\documentclass[12pt]{article}
\usepackage[utf8]{inputenc}
\usepackage[T1]{fontenc}
\usepackage{hyperref}
\title{Wiradjuri Language Resources for K-12 Curriculum Development}
\author{Compiled by Claude Sonnet 5.5 (claude-sonnet-5-5) for Andr\'as Kornai}
\date{}
\begin{document}
\maketitle

\section{Introduction}
Wiradjuri (ISO 639-3 \texttt{wrh}, Glottocode wira1262) is a Pama-Nyungan language of central New South Wales, Australia. It is being revived, with community classes, preschool and school teaching, and published workbooks and dictionaries. The modern Wiradjuri used in many resources is a reconstruction based on historical records.

\section{Primary Educational Resources}
\begin{itemize}
\item Learn Wiradjuri: Volume 1 Ngumbaay (SBS Learn), workbooks by Lynette Riley and Diane Riley-McNaboe, aligned with the Australian Curriculum: \url{https://www.sbs.com.au/learn/?p=4712}
\item The same workbook series has a Volume 2 covering family terms, body parts and colours. No separate URL for Volume 2 was seen.
\item Turvey Park Public School students made a Wiradjuri picture book and dictionary, described in a Daily Advertiser article: \url{https://www.dailyadvertiser.com.au/story/6007396/primary-students-launch-their-first-wiradjuri-dictionary/}
\item Wiradjuri Aboriginal preschool in Griffith teaches the language (Area News): \url{https://www.areanews.com.au/story/5506694/wiradjuri-aboriginal-preschool-teaches-wiradjuri-language/}
\item Area News report on language and identity: \url{https://www.areanews.com.au/story/7141412/language-is-important-for-identity/}
\item Article on teaching children Wiradjuri through games and activities: \url{https://aussiechildcarenetwork.com.au/articles/teaching-children/teach-children-the-aboriginal-language-of-wiradjuri-through-games-and-activities}
\end{itemize}

\section{Linguistic and Reference Materials}
\begin{itemize}
\item AIATSIS AUSTLANG entry D10 Wiradjuri: \url{https://collection.aiatsis.gov.au/austlang/language/D10}
\item Stan Grant Snr and John Rudder, \emph{A New Wiradjuri Dictionary} (2005), is the dictionary the workbooks draw on. No URL for it was seen in the search results.
\end{itemize}

\section{Religious and Community Resources}
\begin{itemize}
\item The Parkes Wiradjuri Language Group produces resources and runs community classes. Parkes Phoenix, ``Spread the word in Wiradjuri'': \url{https://www.parkesphoenix.com.au/spread-the-word-in-wiradjuri/?amp=1}
\item Forbes Phoenix, ``Sharing the Wiradjuri language'': \url{https://www.forbesphoenix.com.au/sharing-the-wiradjuri-language/?amp=1}
\item The Galari-Lachlan Reconciliation Group supplied a school package of the dictionary and DVDs to 14 institutions, described in the Parkes Phoenix article above.
\end{itemize}

\section{Digital and Audio Resources}
\begin{itemize}
\item SBS NITV, ``The two aunties building an understanding of Wiradjuri culture'': \url{https://www.sbs.com.au/nitv/article/the-two-aunties-building-an-understanding-of-wiradjuri-culture/elimw65cv}
\item The DVD ``Yindyamarra Yambuwan'' (15 minutes, on respecting everything) and a children's songs DVD are mentioned in news coverage. No direct URL was seen. A Wiradjuri dictionary app is also mentioned in the Area News preschool article.
\end{itemize}

\section{Curriculum Development Context}
Wiradjuri has the most classroom-ready material of the four languages in this batch. This includes curriculum-aligned workbooks, a dictionary, and local school and preschool programmes. A curriculum should be developed with Wiradjuri community language groups. The AUSTLANG entry dates from about 2012, and the availability of individual items should be checked before relying on them.
\end{document}